% Fixed129 independent substantive insertion.
% A strict integral natural-coefficient marked Jacobian contraction,
% Schreier evaluation projector and derived Hecke equivalence
% on the genuine non-pro-2 S3 residual sector of fixed128.

\subsection{Strict natural-coefficient descent in the non-pro-$2$ $S_3$ sector}
\label{subsec:dyadic-s3-natural-strict-transfer}

The coefficient-specific adjoint Smith certificate of
Theorem~\ref{thm:dyadic-s3-marked-adjoint-jacobian}
does not, by itself, give an explicit map to the
Sylow-preimage cochains.  We now obtain a stronger,
\emph{strict chain-level} result for the distinct
natural rank-two coefficient lattice $T$.
Unlike the adjoint module, the natural $S_3$ module
satisfies the integral norm-zero relation
$I+U_0+U_0^2=0$.  This kills the pro-$2$
projection of every affine lift of the tame
order-three generator, and is precisely the
extra algebra needed for an integral closed
formula at all chain degrees.

Continue with $\Gamma=G_{\mathbf Q_2}$,
$H=G_{\mathbf Q_2(\sqrt[3]{2})}$,
$T=\mathbf Z_2^2$, and the fixed matrices
\[
S=S_0=\begin{pmatrix}0&1\\1&0\end{pmatrix},
\qquad
U=U_0=\begin{pmatrix}-1&-1\\1&0\end{pmatrix}
\]
of \eqref{eq:dyadic-s3-integral-matrices}.
Write $I=I_2$ and put
\begin{equation}\label{eq:dyadic-s3-natural-tame-blocks}
\begin{aligned}
R_\sigma&=S-I=
\begin{pmatrix}-1&1\\1&-1\end{pmatrix},
&
R_\tau&=U-I=
\begin{pmatrix}-2&-1\\1&-1\end{pmatrix},
\\
L_\sigma&=S^{-1}(U-I)
=\begin{pmatrix}1&-1\\-2&-1\end{pmatrix},
&
L_\tau&=S^{-1}-I-U
=\begin{pmatrix}0&2\\0&-1\end{pmatrix}.
\end{aligned}
\end{equation}
The integral determinants are
\begin{equation}\label{eq:dyadic-s3-natural-unit-dets}
\det(R_\tau)=3,\qquad
\det(L_\sigma)=-3,\qquad
\det(S)=-1.
\end{equation}
All these are units in $\mathbf Z_2$.

\begin{theorem}[Exact natural-coefficient full-marked Fox matrix]
\label{thm:dyadic-s3-natural-full-marked-matrix}
Let
$\mathcal J^\bullet_{\Gamma,\mathrm{nat}}(T)$
be the linear \emph{coefficient}
crossed-derivation complex obtained
from the full Roe--Turturean
marked relations in the module $T$
(not the adjoint tangent module
$\operatorname{ad}T$):
\begin{equation}\label{eq:dyadic-s3-natural-three-term-complex}
T\xrightarrow{d^0}T^{\oplus4}
\xrightarrow{d^1}T^{\oplus2},
\end{equation}
with degree-one source coordinates
$(a,b,c,d)$, meaning values at
$(\sigma,\tau,x_0,x_1)$, and target
coordinates in the (tame,wild)
relation order.  The \emph{complete
integral} matrices are
\begin{equation}\label{eq:dyadic-s3-natural-differentials}
\boxed{
\begin{aligned}
d^0(v)&=(R_\sigma v,\ R_\tau v,\ 0,\ 0),
\\
d^1(a,b,c,d)
&=(L_\sigma a+L_\tau b,\ -2c+S^{-1}d).
\end{aligned}}
\end{equation}
Equivalently, in the ordinary
coordinate basis the $4\times8$
degree-one matrix is
\begin{equation}\label{eq:dyadic-s3-natural-4x8-jacobian}
\boxed{\quad
d^1=
\begin{pmatrix}
 1&-1&0&2&0&0&0&0\\
 -2&-1&0&-1&0&0&0&0\\
 0&0&0&0&-2&0&0&1\\
 0&0&0&0&0&-2&1&0
\end{pmatrix}.
\quad}
\end{equation}
Both the tame block $L_\sigma$
and the wild $x_1$ block
$S^{-1}$ are invertible
integrally, while the $\tau$
component of $d^0$ is
$R_\tau$, another integral unit.
Thus $d^1$ has four integral
unit elementary divisors,
and $d^0$ has two further
unit pivots.  After cancelling
these four degree-one--two and
two degree-zero--one contractible
summands, its remaining degree-one
cohomology is free of rank two:
\begin{equation}\label{eq:dyadic-s3-natural-marked-cohomology}
H^0(\mathcal J_{\Gamma,\mathrm{nat}})=0,
\qquad
H^1(\mathcal J_{\Gamma,\mathrm{nat}})
\cong T,\qquad
H^2(\mathcal J_{\Gamma,\mathrm{nat}})=0.
\end{equation}
These are the genuine integral
$G_{\mathbf Q_2}$ cohomology groups
of $T$ in
Theorem~\ref{thm:dyadic-s3-integral-sector}.
\end{theorem}

\begin{proof}
For the coefficient module $T$,
a crossed derivation is a section
of the affine semidirect group
$T\rtimes_\rho S_3$, whose elements
are written $(A,v)$ with law
\[
(A,v)(B,w)=(AB,\ v+Aw).
\]
On the four marked generators
take $(S,a),(U,b),(I,c),(I,d)$.
The tame relation is
$r_t=\sigma^{-1}\tau\sigma\tau^{-2}$;
the affine calculation gives
\[
c(r_t)
=S^{-1}(U-I)a+(S^{-1}-I-U)b,
\]
exactly the first row block of
\eqref{eq:dyadic-s3-natural-differentials}.

The key cancellation is the exact
integral identity
\[
I+U+U^2=0.
\]
For every $v\in T$,
$(U,v)^3=(I,(I+U+U^2)v)=(I,0)$.
Thus every affine $2$-primary
projection $(U,v)^{\omega_2}$
is exactly the identity, while
the projection of $(S,a)$ is
$(S,a)$ because its image
is pro-$2$.
Consequently, in the literal
Roe--Turturean auxiliary words
\eqref{eq:dyadic-omega2-binomial-power}
and Theorem~\ref{thm:dyadic-q2-marked-word-jet-solver},
\[
\sigma_2=\sigma,\quad
u_0=u_1=1,\quad
d_0=x_0^{-1},\quad c_0=h_c=1.
\]
The wild generators act by pure
translations; those translations
commute, and $g_0=\sigma^2$
acts trivially on the translation
subgroup because its linear part
is $S^2=I$.
It follows exactly, not only
modulo $2$, that
\[
h_0=(I,-2c),\qquad
x_1^\sigma=(I,S^{-1}d),
\]
and hence
$c(r_{\rm wild})=-2c+S^{-1}d$.
This proves the full wild row.

The expression for $d^0$
is the principal crossed
derivation $(\rho(g)-I)v$.
The group relator identity implies
$d^1d^0=0$; it also follows
directly from
\eqref{eq:dyadic-s3-natural-tame-blocks}.
By \eqref{eq:dyadic-s3-natural-unit-dets}
the tame $L_\sigma$ and wild
$S^{-1}$ blocks give four
independent unit pivots in
$d^1$, hence $\operatorname{coker}d^1=0$.
The $\tau$ block $R_\tau$
of $d^0$ has two independent
unit pivots, so the image
of $d^0$ is a saturated rank-two
submodule of $\ker d^1$.
Thus $H^0=H^2=0$ and
$H^1$ is free of rank
$8-4-2=2$, as claimed.
The identification with intrinsic
cohomology is made explicit in
the next theorem rather than
assuming the marked relator
complex is a full universal
group-algebra resolution.
\end{proof}

\begin{theorem}[A strict integral deformation retract of the full marked complex]
\label{thm:dyadic-s3-natural-strict-retract}
Define maps of graded
$\mathbf Z_2$-modules between
$\mathcal J^\bullet_{\Gamma,\mathrm{nat}}(T)$
and the single finite module
$T[-1]$:
\begin{equation}\label{eq:dyadic-s3-natural-pi-maps}
\begin{aligned}
p^1(a,b,c,d)&=c,
 &p^0=p^2&=0,\\
i^1(t)&=(0,0,t,2St),
 &i^0=i^2&=0.
\end{aligned}
\end{equation}
Define degree-minus-one
maps $h^1:T^{\oplus4}\to T$ and
$h^2:T^{\oplus2}\to T^{\oplus4}$
by the \emph{closed integral}
matrices
\begin{equation}\label{eq:dyadic-s3-natural-explicit-homotopy}
\boxed{
\begin{aligned}
h^1(a,b,c,d)&=(U-I)^{-1}b,
\\
h^2(u,w)&=(L_\sigma^{-1}u,\ 0,\ 0,\ Sw).
\end{aligned}}
\end{equation}
Then $i,p$ are strict chain
maps and
\begin{equation}\label{eq:dyadic-s3-natural-contraction-identities}
\boxed{\qquad
p i=I_T,\qquad
i p=I_{\mathcal J}-(d h+h d),
\qquad
h^1d^0=I_T,\quad d^1h^2=I_{T^2}.
\qquad}
\end{equation}
In particular
\begin{equation}\label{eq:dyadic-s3-natural-complex-retract}
\mathcal J^\bullet_{\Gamma,\mathrm{nat}}(T)
\ \simeq_{\mathrm{chain},\,\mathbf Z_2}\
T[-1]
\end{equation}
is a \emph{strict, explicit
integral chain-homotopy equivalence},
not just a derived equivalence
from coincident cohomology ranks.

The class $t\in T$ corresponds
to the unique normalized continuous
$\Gamma$-cocycle $c_t$ satisfying
\begin{equation}\label{eq:dyadic-s3-natural-normalized-cocycle}
\boxed{
c_t(\sigma)=c_t(\tau)=0,\qquad
c_t(x_0)=t,\qquad
c_t(x_1)=2St.
}
\end{equation}
Every continuous $\Gamma$-cocycle
with coefficients in $T$ is
cohomologous to exactly one
$c_t$.  Thus evaluation at the
wild source generator $x_0$
gives an intrinsic-to-the-marked-source
isomorphism
\begin{equation}\label{eq:dyadic-s3-natural-cohom-coordinate}
H^1(\Gamma,T)\xrightarrow{\;\sim\;}T,
\qquad[c]\longmapsto c(x_0),
\end{equation}
which is integral and requires
no Fox-to-Herr parametrix.
\end{theorem}

\begin{proof}
The $x_0$ action on $T$ is
trivial, so $p^1d^0=0$.
Equation
\eqref{eq:dyadic-s3-natural-differentials}
gives $d^1i^1=0$, while
$p^1i^1=I_T$.
Hence $i,p$ are chain maps
and one-sided inverses.

The tame block identity
\[
L_\sigma R_\sigma+
L_\tau R_\tau=0
\]
follows from $d^1d^0=0$.
Both $L_\sigma$ and $R_\tau$
are invertible over $\mathbf Z_2$,
so
\[
R_\sigma R_\tau^{-1}
=-L_\sigma^{-1}L_\tau.
\]
Using this identity and
\eqref{eq:dyadic-s3-natural-differentials},
one checks for every
$(a,b,c,d)\in T^4$:
\[
\begin{aligned}
d^0h^1(a,b,c,d)&=
(R_\sigma R_\tau^{-1}b,\ b,\ 0,\ 0),\\
h^2d^1(a,b,c,d)&=
(a+L_\sigma^{-1}L_\tau b,\
 0,\ 0,\ d-2Sc).
\end{aligned}
\]
Their sum is exactly
$(a,b,c,d)-i p(a,b,c,d)$.
Likewise $h^1d^0=I_T$ and
$d^1h^2=I_{T^2}$ because
$L_\sigma^{-1}$, $(U-I)^{-1}$,
and $S$ have integral entries.
This proves all identities in
\eqref{eq:dyadic-s3-natural-contraction-identities}.

The generator values in
\eqref{eq:dyadic-s3-natural-normalized-cocycle}
satisfy both marked word
relations by the exact affine
calculation in
Theorem~\ref{thm:dyadic-s3-natural-full-marked-matrix}.
In each finite quotient
$T/2^nT\rtimes S_3$, the normal
closure of the images of the
wild generators lies in the
$2$-group $T/2^nT$;
the pro-$2$ marked condition is
therefore satisfied.  The
Roe--Turturean presentation
\cite{RoeTurtureanGQ2} supplies
a compatible family of
finite representations, whose
inverse limit defines the
continuous cocycle $c_t$.
Conversely, every continuous
cocycle determines a generator
tuple satisfying the linearized
marked relations, and the
contraction above subtracts
a unique principal cocycle
to reach the displayed
normalized tuple.
As $x_0$ acts trivially,
$c(x_0)$ is unchanged by a
principal coboundary.
This proves the uniqueness
and the actual cohomology
isomorphism.
\end{proof}

\begin{theorem}[Explicit index-three Sylow cocycle projector]
\label{thm:dyadic-s3-explicit-cocycle-projector}
Let
$t_i=\tau^i$ for $i=0,1,2$
be the right-coset representatives
of Theorem~\ref{thm:dyadic-s3-schreier-corestriction},
and write
\[
w_{ij}=\tau^i x_j\tau^{-i}\in H
\qquad(i=0,1,2;\ j=0,1).
\]
For \emph{every} continuous
$H$-cocycle $c\in Z^1(H,T)$
define the finite integral
Schreier-coordinate functional
\begin{equation}\label{eq:dyadic-s3-strict-phi-functional}
\boxed{\quad
\Phi_H(c)=\frac13
\sum_{i=0}^{2}U^{-i}c(w_{i0})\in T.
\quad}
\end{equation}
This descends to an
integral surjection on
$H^1(H,T)$ because each
$w_{i0}$ acts trivially
on the coefficient lattice.
The normalized global cocycles
of \eqref{eq:dyadic-s3-natural-normalized-cocycle}
satisfy
\begin{equation}\label{eq:dyadic-s3-strict-res-values}
\boxed{\quad
\begin{gathered}
(\operatorname{res}c_t)(\sigma)=0,\qquad
(\operatorname{res}c_t)(\tau^3)=0,\\
(\operatorname{res}c_t)(w_{i0})=U^i t,
\qquad
(\operatorname{res}c_t)(w_{i1})=2U^iS t,
\quad i=0,1,2.
\end{gathered}\quad}
\end{equation}
The restriction map is
injective in degree one,
and its left inverse is
$\Phi_H$:
\begin{equation}\label{eq:dyadic-s3-strict-res-phi}
\Phi_H(\operatorname{res}c_t)=t.
\end{equation}
More strongly, for every
$H$-\emph{cocycle}, not just
its cohomology class,
\begin{equation}\label{eq:dyadic-s3-strict-e-cocycle}
\boxed{\quad
\frac13\operatorname{res}
\operatorname{cor}(c)
=\operatorname{res}c_{\Phi_H(c)}
\qquad\text{as literal }H
\text{-cocycles}.\quad}
\end{equation}

In the eight displayed
Schreier generator coordinates
\[
(\sigma,\tau^3,\
w_{00},w_{10},w_{20},\
w_{01},w_{11},w_{21}),
\]
the following \emph{entirely
finite} block matrices describe
the degree-one projector:
\begin{equation}\label{eq:dyadic-s3-schreier-16x16-projector}
\boxed{
\begin{aligned}
\mathsf I_{\rm Sch}&=
\begin{pmatrix}
0\\0\\I\\U\\U^2\\2S\\2US\\2U^2S
\end{pmatrix}:T\longrightarrow T^8,\\[3pt]
\mathsf P_{\rm Sch}&=\frac13
\begin{pmatrix}
0&0&I&U^{-1}&U^{-2}&0&0&0
\end{pmatrix}:T^8\longrightarrow T,\\
\mathsf E_{\rm Sch}
&=\mathsf I_{\rm Sch}\mathsf P_{\rm Sch},
\qquad
\mathsf E_{\rm Sch}^2=\mathsf E_{\rm Sch}.
\end{aligned}}
\end{equation}
Every block is an integral
$2\times2$ matrix; the resulting
$\mathsf E_{\rm Sch}$ is a
strict rank-two $16\times16$
projection on the ambient
Schreier evaluation space.
On actual cocycles it is
\emph{exactly} the normalized
restriction--corestriction
operator.  On cohomology
\begin{equation}\label{eq:dyadic-s3-h1-primitive-decomposition}
\boxed{\quad
H^1(H,T)
=\operatorname{res}H^1(\Gamma,T)
\oplus\ker\Phi_H
\cong\mathbf Z_2^2
\oplus\mathbf Z_2^5.
\quad}
\end{equation}
The formula does not presume
that these eight displayed
Schreier words form a
\emph{minimal} pro-$2$
Demu\v{s}kin presentation.
\end{theorem}

\begin{proof}
Every $w_{ij}$ lies in
the wild normal subgroup
of the marked Galois
presentation and maps
to the identity under $\rho_T$.
Thus $\Phi_H$ kills
$H$-coboundaries.
The crossed-cocycle formula
and $\rho_T(\tau)=U$ give
\[
c_t(\tau^i x_j\tau^{-i})
=U^i c_t(x_j)
\]
because $c_t(\tau)=0$ and
$\rho_T(x_j)=I$.
This gives
\eqref{eq:dyadic-s3-strict-res-values}
and
\eqref{eq:dyadic-s3-strict-res-phi}.

For arbitrary $c\in Z^1(H,T)$,
the element $\tau^3\in H$
has pro-odd order and
acts trivially on $T$.
Continuity forces
$c(\tau^3)=0$.
The coset formula
\eqref{eq:dyadic-s3-transfer-one}
then gives
\[
(\operatorname{cor}c)(\sigma)
=(I+U^{-1}+U^{-2})c(\sigma)=0,
\quad
(\operatorname{cor}c)(\tau)=0,
\]
and
\[
(\operatorname{cor}c)(x_0)
=\sum_{i=0}^{2}U^{-i}c(w_{i0})
=3\Phi_H(c).
\]
By the exact normalized-cocycle
description in
Theorem~\ref{thm:dyadic-s3-natural-strict-retract},
a $\Gamma$-cocycle with
zero tame-generator values
is determined uniquely by
its value at $x_0$ and
is $c_t$ with that value.
Thus $\operatorname{cor}c
=c_{3\Phi_H(c)}$,
and dividing by the odd
unit $3$ gives the exact
cocycle identity
\eqref{eq:dyadic-s3-strict-e-cocycle}.

Matrix multiplication shows
\[
\mathsf P_{\rm Sch}
\mathsf I_{\rm Sch}
=\tfrac13(I+U^{-1}U+U^{-2}U^2)
=I,
\]
hence $\mathsf E_{\rm Sch}$
is a strict integral
idempotent of rank two.
For actual $H$-cocycles,
the matrix is the evaluation
of \eqref{eq:dyadic-s3-strict-e-cocycle}.
The rank-seven freeness of
$H^1(H,T)$ from
Theorem~\ref{thm:dyadic-s3-integral-sector}
now gives the rank-five kernel
and direct sum
\eqref{eq:dyadic-s3-h1-primitive-decomposition}.
\end{proof}

\begin{theorem}[A strict integral bar-cochain Hecke representative]
\label{thm:dyadic-s3-strict-bar-hecke-projector}
Let $C^\bullet(H,T)$ be the
normalized continuous $H$-bar
cochain complex.  Extend
\eqref{eq:dyadic-s3-strict-phi-functional}
to \emph{arbitrary} continuous
degree-one cochains
$f:H\to T$, and set
\[
p_H^1(f)=\frac13\sum_{i=0}^{2}
U^{-i}f(w_{i0}),\quad
p_H^k=0\ (k\ne1).
\]
Define the chain map
$i_H:T[-1]\to C^\bullet(H,T)$
by $i_H^1(t)=\operatorname{res}c_t$,
and $i_H^k=0$ for $k\ne1$.
Then $p_H:C^\bullet(H,T)\to T[-1]$
is a \emph{strict cochain map},
$p_Hi_H=I_T$, and
\begin{equation}\label{eq:dyadic-s3-strict-bar-projector}
\boxed{\quad
\Pi_H=i_Hp_H,\qquad
\Pi_H^2=\Pi_H,\qquad
\mathsf T_{\rm Hk}^{\rm str}=3\Pi_H-I.
\quad}
\end{equation}
Both the idempotent identity
and the Hecke polynomial hold
as \emph{literal identities of
integral cochain maps}:
\begin{equation}\label{eq:dyadic-s3-strict-bar-hecke-polynomial}
\boxed{\quad
(\mathsf T_{\rm Hk}^{\rm str})^2
=\mathsf T_{\rm Hk}^{\rm str}+2I.
\quad}
\end{equation}
The strict operator
$\Pi_H$ represents the
\emph{canonical derived}
Sylow idempotent:
\begin{equation}\label{eq:dyadic-s3-bar-derived-equality}
[\Pi_H]=
\left[
\tfrac13\operatorname{res}_H^\Gamma
\operatorname{cor}_H^\Gamma
\right]
\quad
\text{in }
\operatorname{End}_{D(\mathbf Z_2)}
R\Gamma(H,T).
\end{equation}

There is also an explicit
integral chain-level comparison
between the \emph{finite}
full-marked coefficient
complex and the \emph{image}
of the Sylow bar-cochain
projector.  With $i,p$ from
Theorem~\ref{thm:dyadic-s3-natural-strict-retract},
set
\begin{equation}\label{eq:dyadic-s3-marked-to-sylow-chain-maps}
\mathsf F=i_Hp:
\mathcal J^\bullet_{\Gamma,\mathrm{nat}}
\longrightarrow C^\bullet(H,T),
\qquad
\mathsf G=i p_H:
C^\bullet(H,T)
\longrightarrow\mathcal J^\bullet_{\Gamma,\mathrm{nat}}.
\end{equation}
These are \emph{strict integral}
cochain maps and satisfy
\begin{equation}\label{eq:dyadic-s3-marked-to-sylow-composites}
\boxed{\quad
\mathsf G\mathsf F=i p
\simeq I_{\mathcal J}
\ \text{by the explicit }h,
\qquad
\mathsf F\mathsf G=\Pi_H
\ \text{strictly}.
\quad}
\end{equation}
In particular the entire
degree-one Hecke/Sylow
summand is finite and
directly reconstructed from
the marked $x_0$-coordinate.
This closes the natural-$T$
\emph{integral descent chain}
without selecting a
five-generator Demu\v{s}kin
Fox basis for $P_H$.

The statement does \emph{not}
upgrade the actual infinite
bar complex to a globally
finite free group-algebra
resolution.  It identifies
a \emph{strict finite
cochain summand} of it.
\end{theorem}

\begin{proof}
For a $0$-cochain $v\in T$,
the continuous bar coboundary
has
\[
(d^0v)(w_{i0})
=(\rho_T(w_{i0})-I)v=0.
\]
Thus $p_H^1d^0=0$,
and all other cochain
map conditions for
$p_H:C^\bullet(H,T)\to T[-1]$
are automatic.
The chosen $i_H^1(t)$
is a continuous cocycle,
so $i_H$ is a chain map.
Equation
\eqref{eq:dyadic-s3-strict-res-phi}
gives $p_Hi_H=I_T$.
Consequently $\Pi_H$
is strictly idempotent.
Expanding
$(3\Pi_H-I)^2$
using $\Pi_H^2=\Pi_H$
gives
\eqref{eq:dyadic-s3-strict-bar-hecke-polynomial}.

To compare with the canonical
derived transfer, define
$p_\Gamma:C^\bullet(\Gamma,T)\to T[-1]$
by $p_\Gamma^1(f)=f(x_0)$
and zero in other degrees.
It is a cochain map because
$x_0$ acts trivially on $T$.
Let $i_\Gamma^1(t)=c_t$.
Theorem~\ref{thm:dyadic-s3-natural-strict-retract}
proves $p_\Gamma i_\Gamma=I$
and that $p_\Gamma$ and
$i_\Gamma$ induce inverse
isomorphisms on intrinsic
$G_{\mathbf Q_2}$ cohomology.
Hence their composites are
identities in $D(\mathbf Z_2)$.

The explicit degree-one
corestriction formula
\eqref{eq:dyadic-s3-transfer-one}
holds for arbitrary
continuous cochains as
well as cocycles, and
gives the \emph{strict}
identity of cochain maps
\[
p_H=\tfrac13 p_\Gamma
\,\operatorname{cor}_H^\Gamma,
\qquad
i_H=\operatorname{res}_H^\Gamma
\,i_\Gamma.
\]
Consequently
\[
[\Pi_H]=
\left[\operatorname{res}\,
i_\Gamma\,\tfrac13p_\Gamma
\operatorname{cor}\right]
=\left[\tfrac13\operatorname{res}
\operatorname{cor}\right]
\]
in the derived category,
proving
\eqref{eq:dyadic-s3-bar-derived-equality}.
Finally all maps in
\eqref{eq:dyadic-s3-marked-to-sylow-chain-maps}
are composites of strict
cochain maps, and
$p_Hi_H=I$, $pi=I$
imply
\[
\mathsf G\mathsf F
=i(p_Hi_H)p=ip,\qquad
\mathsf F\mathsf G
=i_H(pi)p_H=i_Hp_H.
\]
The homotopy in
\eqref{eq:dyadic-s3-natural-explicit-homotopy}
satisfies $ip\simeq I_{\mathcal J}$
over $\mathbf Z_2$,
giving the claimed chain-level
comparison.
\end{proof}

\begin{warning}[What is and is not native in fixed129]
\label{warn:dyadic-s3-natural-minimal-fox-boundary}
The natural-coefficient descent
projector is now strict on
the integral bar complex,
is an actual finite
$16\times16$ block matrix
on the displayed Schreier
generator-evaluation space,
and is a strict chain
summand equivalent to the
finite full-marked
coefficient Jacobian.
Its image is the same
derived object as the
all-Sylow Fox retract.
Nevertheless an
\emph{explicit identification
of the eight Schreier words
with a specific minimal
five-generator Demu\v{s}kin
basis}, including the
integer relation module,
its Fox syzygies and the
direct finite
Schreier-to-minimal-Fox
second-degree chain map,
has not been constructed.
Nor does the norm-zero
identity hold for the
adjoint $S_3$ module:
$\operatorname{ad}T$ has
nonzero tame $C_3$ invariants
and requires the separate
eight-by-sixteen Smith
certificate of fixed128.
The theorem is therefore
a sharp result for the
fixed natural coefficient
carrier, not an assertion
of an all-representation
full-group Fox resolution.
\end{warning}
